\documentclass{sapientia-thesis}
\begin{document}

\chapter{Szakdolgozat Írási és Formázási Útmutató}
Ez a dokumentum a Sapientia egyetemi sablon eredeti útmutatóit tartalmazza segédletként.

\section{Szakdolgozat fejezetek és tartalmi követelmények útmutatója}
Nagyon nehéz általános szabályokat felállítani, vagy kőbe vésett arányokat, oldalszámokat mondani, hiszen minden munka egy kicsit más. Itt mégis megpróbálom felvázolni nektek egy általános diplomaterv/szakdolgozat szerkezetét, amitől természetesen el lehet térni, amennyiben a téma ezt megkívánja.

\subsection{Fejezetek}
A dolgozat számos fejezetbõl épül fel. Fontos látni, hogy ezeket egyáltalán nem szükséges sorrendben feltölteni. A bevált gyakorlat az, hogy e sablonból kiindulva kitöröljük az ismertetõ részeket, \verb|\chapter|--\verb|\subsection| parancsokkal összeállítjuk a főbb fejezeteket, amikről írni akarunk, kb. olyan terjedelemben, hogy a tartalomjegyzék 1--1.5 oldal hosszú legyen, majd az aktuális munkától függően tetszőleges sorrendben elkezdjük őket tartalommal feltölteni. Ha megakadunk és úgy érezzük, hogy már mindent leírtunk (és még csak 15 oldalnál tartunk), olvassuk el egy évfolyamtárssal egymás dolgozatát és írjunk bele kérdéseket mindenhova, ahol valami nem egyértelmű. De nézzük meg, mik is a nélkülözhetetlen szerkezeti elemek.

\subsection{Feladatkiírás}
Nem sorszámozott oldal, a címoldalt követi a dolgozatban.

\subsection{Címoldal}
A diplomaterv címe, egyetem, tanszék, saját neved, konzulens neve és a védés éve. \LaTeX-ben a \verb|\studentname|, \verb|\titleHU|, \verb|\supervisorHU| stb.\ parancsokkal töltsd ki a sablon elején lévő metaadatokat.

\subsection{Tartalomjegyzék}
A tartalomjegyzéket a \LaTeX\ automatikusan generálja a \verb|\chapter|, \verb|\section| és \verb|\subsection| parancsokból -- nincs szükség kézi frissítésre. A sablon eleve az elsõ 3 szintet jeleníti meg (\verb|\setcounter{tocdepth}{3}|). A tartalomjegyzék pontosan akkor frissül, amikor újrafordítod a dokumentumot; ha hivatkozások is változtak, futtasd le a teljes fordítási sorozatot: \texttt{pdflatex} $\to$ \texttt{biber} $\to$ \texttt{pdflatex} $\to$ \texttt{pdflatex}.

\subsection{Nyilatkozat}
A törlendõ részeket húzd ki, a neved, dátumot töltsd ki mindenhol. Ne felejtsd el aláírni sem.

\subsection{Bevezetés}
Itt kezdődik a dolgozat érdemi része. A bevezetés célja a feladat értelmezése, a motiváció leírása, a kontextus megteremtése. Ami a stílust illet, képzeld azt, hogy egy kívülállónak próbálod elmagyarázni, mit csináltál és miért. Indíts messzirõl, pl. beszélj a széles értelemben vett területed fontosságáról, majd fokozatosan közelíts rá arra a szűkebb problémára amit látsz, s amit megpróbálsz megoldani a dolgozatodban.

Tipikus hiba szokott lenni, hogy motivációként azt írod le, hogy „a konzulensemtõl ezt a feladatot kaptam". Ez nagyon bénán hangzik. Próbáld inkább egy tágabb perspektívából nézni a dolgot. A bevezetés praktikus zárása egy olyan alfejezet, ami a dolgozat további szerkezetét ismerteti.

\subsection{Irodalomkutatás, technológiák, hasonló alkotások bemutatása}
Ennek a fejezetnek a lényege, hogy bemutassa azt az alapot, amire a munkádat építetted. Itt bemutathatod a szakirodalmat, mások munkáját, eredményeit, amikhez képest a sajátodat majd később meg tudod határozni. Fontos, hogy mindig a témád szempontjából releváns részekről írj. Ez a fejezet azoknak szól, akik szakmabeliek, de nem konkrétan ezzel a területtel foglalkoznak.

\subsection{A feladatkiírás pontosítása és részletes elemzése}
Tipikus fejezetcím lehet az „Architektúra", vagy „Tervezés". A feladat itt már egyáltalán nem mesélgetés, hanem egy objektív elemzés arról, hogy mi is a konkret feladat, milyen funkcionális és nem funkcionális követelmények vannak, hogy néz ki a rendszer architektúrája amit terveztél.

\subsection{Önálló munka bemutatása}
Ez a blokk egy, vagy több nagyobb fejezetet tartalmaz, mely bemutatja az elkészült megoldás műszaki felépítését, kiemelve az érdekesebb/bonyolultabb megoldásokat és a nem egyértelmű technológiai döntések indoklásait.

Nagyon fontos, hogy szemben az előző blokkal itt nem cél, hogy egy száraz dokumentációt készítsél. Tipikus hiba tud lenni, hogy minden apró részletet és függvényt be akarsz mutatni. Tervezz meg egy gondolati szálat, ami mentén be akarod mutatni a rendszert, és azon belül se mindent próbálj bemutatni, hanem koncentrálj az érdekesebb, nem triviálisabb megoldásokra.

\subsection{Önálló munka értékelése, mérések, tesztelés, eredmények bemutatása}
E fejezet célja, hogy összegezve bemutassa az eredményeidet. Például ha valamilyen algoritmust fejlesztettél, itt mutathatod be mérésekkel, diagramokkal, hogy milyen teljesítményt produkál. Ha egy weboldalt vagy mobil alkalmazást készítettél, és készülttek hozzá tesztek, azokat itt bemutathatod.

\subsection{Összefoglaló}
Dolgozattípustól függetlenül 1 oldalban összefoglalóként írd le az eredményeidet. E/1-ben és múlt időben. Megterveztem, megvalósítottam, eldöntöttem, leteszteltem\ldots stb.

\subsection{Köszönetnyilvánítások}
Őszintén szólva, a legtöbb esetben felesleges blokk, pusztán a konzulensednek megköszönni a segítségét nem szükséges. Akkor érdemes ezzel foglalkozni, ha valamilyen harmadik fél is érdemben segítette a munkádat.

\subsection{Részletes és pontos irodalomjegyzék}
Ezt a \texttt{biblatex}/\texttt{biber} páros generálja neked automatikusan a \verb|\printbibliography| parancs hatására.

\subsection{Függelék}
A függelékek a törzstartalmon kívüli kiegészítések a dolgozathoz. Ide kerülhetnek a nagyobb ábrák, hosszabb példakódok, amelyek csak lazábban kapcsolódnak a témához. A függelék nem kötelező tartalmi elem.

\subsection{Egyéb tartalmi elemek és stílus}
A szakdolgozat/diplomaterv műfaját tekintve félúton van egy szoftverspecifikáció és egy irodalmi esszé között. Fontos, hogy a leírtak objektív elemzés eredményei legyenek, mérnöki szemléletmódot tükrözzenek. Ugyanakkor kerüld a száraz felsorolásokat, s próbáld inkább elbeszélő stílusban bemutatni a munkádat.

A bíráló számára lényeges, hogy határozottan szét tudja választani az önálló munkádat mások munkájától. Ezért rendkívül fontos, hogy következetesen és sűrűn alkalmazd az egyes szám, első személyű mondatokat. Ne azt írd, hogy „Ez itt az osztálydiagram", hanem azt, hogy „A~\ref{fig:pelda_guide}. ábrán bemutatom az osztálydiagramot, amit az alkalmazásomhoz terveztem".

A dolgozatodban alapszabály, hogy \textbf{minden} rövidítést legalább az első használatkor ki kell fejteni. Például: „A weboldal leírására HyperText Markup Language-t (HTML) használtam." -- ezt követően a \texttt{HTML} rövidítés már szabadon használható.

Számos esetben előfordul, hogy egy-egy probléma megoldására több különböző technológia áll rendelkezésedre. Ha a feladatkiírás nem kötötte ki a technológiát, és nem a nyilvánvalóan korszerűbbet választod, indokold meg a döntésed. Mindenképp kerüld viszont az amatőr magyarázkodást és a szubjektív érvelést: „Azért választottam az iOS platformot, mert sokkal jobbnak tartom az Androidnál" -- az ilyen mondatokat előszeretettel szedik cafatokra a bírálók.

\subsection{Mesterséges intelligencia és generatív MI-eszközök}
A Kar mind az oktatókat, mind a hallgatókat támogatja a generatív mesterséges intelligencia (MI) használatában. A generatív MI az elmúlt évtized legnagyobb hatású informatikai innovációja, amely hatékonyabbá teszi a mérnöki munkát. Arra biztatjuk a Kar minden polgárát, hogy az etikus alkalmazás keretein belül maximálisan használja ki a generatív MI képességeit. Az eszközök használata során azonban fontos betartani a Kar irányelveit, kérjük, mindenki figyelmesen tanulmányozza azokat.

\subsection{Dolgozat értékelése}
Szakdolgozat (BSc) esetében a dolgozat értékelése egy kötött pontrendszer szerint történik, mely a következő témaköröket értékeli:
\begin{itemize}
  \item A dolgozat szerkesztése, stílusa
  \item Az adatgyűjtés, irodalomfeldolgozás színvonala
  \item Technológia, eszköz- és módszerválasztás
  \item Kidolgozás színvonala, meghozott döntések indoklása
  \item Kidolgozás mélysége, elvégzett munka mennyisége
  \item A végeredmény (kontribúció) minősége
  \item Az eredmények hallgatói értékelése, validálása
\end{itemize}

\bigskip
\noindent\textbf{Figyelem!} A záróvizsga tervek számára (mérnöki szakok végzettjei esetében) \textbf{50--80 oldalas terjedelem} ajánlott, \textbf{A4 oldalméret}, \textbf{12-es Times New Roman} (vagy annak megfelelő) betűméret, \textbf{1--1,15/1,5 sorköz} + mellékletek.


\section{\LaTeX\ formázási és helyesírási segédlet}
Ez a fejezet összeszedi azokat a nélkülözhetetlen elemeket, amelyeket a dolgozat \LaTeX-ben való elkészítése során használnod kell.

\subsection{Általános tudnivalók}
A diplomaterv szabványos méretű A4-es lapokra kerüljön. Az alapértelmezett betűkészlet a 12 pontos Times New Roman, 1.15-ös sorközzel. Ez a sablon (\texttt{sapientia-thesis.cls}) a fenti előírásoknak megfelelően van kialakítva, így általában külön nem kell foglalkoznod velük.

Minden oldalon -- az első négy szerkezeti elem kivételével -- szerepelnie kell az oldalszámnak.

A fejezeteket decimális beosztással kell ellátni. Az ábrákat a megfelelő helyre be kell illeszteni, fejezetenként decimális számmal és kifejező címmel kell ellátni -- ezt a \LaTeX\ automatikusan kezeli a \verb|\caption| és \verb|\label| parancsok segítségével.

Az irodalomjegyzék szövegközi hivatkozása sorszámozva történik (\texttt{[1]}, \texttt{[2]} stb.). A teljes lista a hivatkozási sorrendben a szöveg végén szerepel. A szakirodalmi források címeit mindig az eredeti nyelven kell megadni. A listában szereplő valamennyi publikációra hivatkozni kell a szövegben.

\subsection{Stílusok}
A \LaTeX\ szövegek egységességét környezetek (\textit{environments}) és parancsok (\textit{commands}) segítségével lehet a legegyszerűbben garantálni. A sablon előre definiált parancsokat tartalmaz a leggyakoribb esetekre (fejezetek, ábrák, kódblokkok, hivatkozások). Azt javaslom, hogy a sablon struktúráját kövesd, és ne definiálj saját formázó parancsokat -- ha valami hiányzik, szólj a konzulensednek.

\subsection{Címsorok}
A fejezetcímek esetén a \verb|\chapter|, \verb|\section|, \verb|\subsection| és \verb|\subsubsection| parancsokat használjuk. Negyedik szintnél mélyebbre egy ilyen terjedelmű munkában ritkán van szükség; ha ez mégis felmerülne, sokszor inkább a fejezetszerkezetet érdemes átgondolni újból.

Kerüljük azt, amikor egy fejezetcím után rögtön alfejezet következik -- legalább pár sor bevezető szöveg legyen közöttük. Hasonlóan érdemes elkerülnünk azt is, amikor a címsor után rögtön felsorolás következik.

\subsection{Másolás, beillesztés}
A copy-paste a szép formázás legnagyobb gyilkosa. Amennyiben külső forrásból (weboldal, PDF, más szerkesztő) illeszted be a szöveget egy \texttt{.tex} fájlba, figyelj arra, hogy az idézőjelek (\verb|„|~és~\verb|"|), kötőjelek (\verb|--|, \verb|---|) és az ékezetes karakterek helyesen kerüljenek át. \LaTeX-ben a speciális karaktereket (\verb|%|, \verb|&|, \verb|$|, \verb|#| stb.) visszaperrel kell escape-elni: \verb|\%|, \verb|\&|, \verb|\$| stb.

Szöveg áthelyezésekor a kereszthivatkozások (\verb|\ref|, \verb|\cite|) automatikusan érvényesek maradnak -- ez a \LaTeX\ egyik nagy előnye a szövegszerkesztőkkel szemben.

\subsection{Fordítás és frissítés}
A \LaTeX-ben a tartalomjegyzék, ábrajegyzék, kereszthivatkozások és irodalomjegyzék automatikusan frissülnek fordításkor -- nincs szükség kézi mezőfrissítésre. Hivatkozások változásakor azonban a teljes fordítási sorozatot le kell futtatni:

\begin{enumerate}
  \item \texttt{pdflatex main} -- első átmenet, aux-fájlok generálása
  \item \texttt{biber main} -- irodalomjegyzék feldolgozása
  \item \texttt{pdflatex main} -- hivatkozások beírása
  \item \texttt{pdflatex main} -- tartalomjegyzék véglegesítése
\end{enumerate}

A legtöbb \LaTeX-szerkesztő (TeXstudio, VS~Code + LaTeX Workshop, Overleaf) ezt a sorozatot egyetlen gombnyomásra elvézi. Ha csak a szöveg változott és nincs új hivatkozás, elegendő egyszer fordítani.

\subsection{Helyesírás}
A rossz helyesírásra nincs mentség. E fejezetben összeszedem a leggyakrabban látott hibákat, amiknek elkerülésére érdemes odafigyelni. Ettől függetlenül melegen ajánlom, hogy a kész dolgozatod olvastasd át egy barátoddal/családtagoddal, hogy az apróbb, megbúvó hibákat is kiszűrd.

\subsubsection{Elgépelések}
Ez mindenkivel megesik. A TeXstudio és a VS~Code~+~LaTeX~Workshop bővítmény beépített helyesírás-ellenőrzővel rendelkezik; győzödj meg arról, hogy a \textbf{magyar szótár} telepítve és beállítva van. Overleaf-en a \textit{Spell check} menüben válaszd a \textit{Hungarian} nyelvet.

\subsubsection{Egyeztetés hiánya}
Az elírások egyik leggyakoribb formája az egyes szám/többes szám egyeztetésének hiánya mondatrészek között, mint például: „Petike és a barátnője elmentek a boltba és hozott egy kiló kenyeret." Ezek a mondatok főleg az utólagos átfogalmazások, belejavítások során keletkeznek; legjobb védelem ellenük az utólagos átolvasás.

\subsubsection{Külföldi szavak, kifejezések}
Az idegen szavakkal csak a baj van. Az általános jó tanácsom, hogy amennyiben csak lehetséges, \textbf{magyar, vagy magyarosított írásmódú szakkifejezéseket használj}, és a könnyebb olvashatóság érdekében mindig kerüld az idegen szavak ragozását. Pl. „Apache-csal" helyett írjuk azt, hogy „Apache webszerverrel".

Néhány gyorstipp szoftverfejlesztőknek: property $\to$ tulajdonság; event $\to$ esemény; method $\to$ metódus/függvény; debug $\to$ hibakeresés; file $\to$ fájl.

\subsubsection{Helyesírás-ellenőrző}
Személyes ízlés kérdése, hogy milyen szerkesztővel dolgozol, ugyanakkor azt \textbf{meg kell oldanod, hogy legyen meg mellé magyar nyelvű helyesírás-ellenőrződ}. E nélkül dokumentumot szerkeszteni olyan, mint papíron programozni.

TeXstudio esetén: \textit{Options / Configure TeXstudio / Editor / Spelling dictionary} alatt telepítse a \texttt{hu\_HU} szótárat. Tesztelésképpen írd be, hogy \textit{„elgemépelés"} -- ha nem húzza alá pirossal, valami nincs rendben a beállításaiddal.

\subsection{Képek}
A dolgozatodban valószínűleg számos ábrára lesz szükséged.

\subsubsection{Beszúrás, formázás}
A képhez használd a \texttt{figure} környezetet, a képaláírást a \verb|\caption| paranccsal add meg, és adj neki egy \verb|\label|-t a hivatkozáshoz. A sablonban az \texttt{[H]} elhelyezési opció (a \texttt{float} csomag) rögzíti az ábrát a szöveghez képest, ami diplomadolgozatnál általában a legjobb választás.

Az ábra sorszáma mellé mindig érdemes rövid magyarázatot is fűzni. Vektorgrafikus ábrákat PDF vagy EPS formátumban illeszd be; ezek nyomtatásban is élesen jelennek meg.

Tapasztalni fogod, hogy a \texttt{figure} és \texttt{table} lebegő elemek néha nem a kívánt helyre kerülnek. Az ábrák méretének módosításával, vagy a \verb|[H]| opció alkalmazásával tudod megoldani, hogy ne maradjanak nagy üres felületek a dolgozatodban.

\subsubsection{Képminőség}
A raszteres (tehát nem vektorgrafikus) képek használata különös körültekintést igényel. Ezek kiválóan néznek ki a 90 dpi-s monitorodon, ám a 600/1200~dpi-s nyomtatókon kinyomtatva rendkívül bénák lesznek a szép, pixelmentes szövegek és vektorgrafikus ábrák mellett.

Ha tehát lehetséges, használjunk vektorgrafikus ábrákat (PDF, EPS, SVG), vagyis a diagramokat, forráskódot stb. ne képernyőképeken keresztül, hanem közvetlen \verb|\includegraphics| megoldással illesszük be. Ha elkerülhetetlen a raszteres képek használata, próbáljunk meg minél magasabb felbontású képet berakni.

\subsection{Kereszthivatkozások}
Amennyiben szeretnél egy ábrára, táblázatra, képletre vagy korábbi fejezetre hivatkozni, helyezz el egy \verb|\label{azonosito}| parancsot a hivatkozott elemnél, majd a szövegben hivatkozz rá a \verb|\ref{azonosito}| paranccsal.

Kerüld az „előző oldalon", „következő oldalon" fordulatokat, ugyanis az ábrák végső helyzete a tördelés során még megváltozhat.

\subsection{Irodalomhivatkozások}
Az irodalomhivatkozások kezelésére a \texttt{biblatex} csomagot és a \texttt{biber} háttérprogramot használjuk. A forrásokat a \texttt{references.bib} fájlban kell felvenni BibTeX formátumban. Amikor dolgozatodban egy külső műre, weboldalra, könyvre, előadásra stb.\ szeretnél hivatkozni, add hozzá a \texttt{.bib} fájlhoz, majd a szövegben hivatkozz rá a \verb|\cite{kulcs}| paranccsal~\cite{piczak2015esc, bosch2012irmas}.

Folyóirat cikkeknél a szerzők mellett szerepeljen a pontos cím, a folyóirat neve, évfolyam, szám és oldalszám. Internet hivatkozásoknál fontos, hogy az oldal szerzőjét, címét és az elérés időpontját is add meg (mivel a link egy idő után elérhetetlenné is válhat) -- a BibTeX \texttt{@online} típusával és az \texttt{urldate} mezővel.

A sablon \texttt{numeric} stílusú hivatkozásokat használ: a szövegben a forrásra \texttt{[1]}, \texttt{[2]} alakban hivatkozunk, a lista a szöveg végén, hivatkozási sorrendben jelenik meg.

\subsubsection{Pozíciónálás}
Az irodalomhivatkozások a szövegtörzsben, ábrák képaláírásaiban és táblázatokban is előfordulhatnak, de \textbf{fejezetcímekben soha}. Ha egy adott forrás egy egész bekezdésre vonatkozik, akkor is elég, ha az első mondat vagy bekezdés után megemlítjük.

\subsubsection{Mikor kell hivatkoznom?}
Minden külső forrásból átvett képnél, szövegrésznél, illetve olyan állítások, technológiák, algoritmusok megemlítésénél, melyek nem feltétlenül egyértelműek egy átlagos műveltségű olvasó számára~\cite{piczak2015esc}. Nagyságrendileg egy szakdolgozatban átlagosan 10--20, egy diplomatervben átlagosan 20--30 külső forrást illik megemlíteni. Ha lehet, törekedjél nyomtatott forrásokra (folyóiratcikkek, könyvek, konferencia-kiadványok), és csak akkor hivatkozz weboldalakra, ha nyomtatott forrást nem találtál.

\subsection{\LaTeX\ tippek és trükkök}

\subsubsection{Navigáció a dokumentumban}
Mivel a fejezeteket nagy valószínűséggel nem sorrendben fogod tartalommal feltölteni, érdemes kihasználni a szerkesztő navigációs lehetőségeit. A TeXstudio bal oldali paneljén a \textit{Structure} nézet mutatja az összes fejezetet és alfejezetet, egy kattintással az adott helyre ugrik. A \texttt{hyperref} csomag (a sablonban már betöltve) a generált PDF-be kattintható könyvjelzőket és belső linkeket szúr be, amelyek PDF-megjelenítőkben navigációs panelként jelennek meg.

\subsubsection{Megjegyzések a forráskódban}
A \LaTeX-ban a megjegyzéseket a \verb|%| karakterrel lehet beilleszteni -- a sor többi részét a fordító figyelmen kívül hagyja. Ezt előszeretettel használhatod, hogy megjelöljed azokat a részeket, ahova még vissza kell térned. Konzulensed a \texttt{changes} csomag segítségével tud nyomtatható megjegyzéseket és nyomonkövethető javításokat fűzni a dokumentumhoz.

\subsubsection{Korrektúra és együttműködés}
Konzulensed legegyszerűbben a \texttt{changes} csomag (\verb|\added{...}|, \verb|\deleted{...}|, \verb|\replaced{...}{...}|) segítségével tud megjegyzéseket és javításokat fűzni a dokumentumhoz, amelyek a PDF-ben jól látható módon jelennek meg. Overleaf esetén a beépített Track Changes funkciót is használhatjátok.

\subsubsection{Hasznos gyorsbillentyűk -- TeXstudio}
A TeXstudio számos beépített gyorsbillentyűt kínál. A legfontosabbak: \texttt{F5} -- fordítás és megjelenítés; \texttt{F6} -- csak fordítás; \texttt{Ctrl+Shift+B} -- \texttt{biber} futtatása; \texttt{Ctrl+Klikk} -- ugrás a forrásból a PDF-be (és vissza). Tetszőleges parancshoz saját gyorsbillentyű rendelhető az \textit{Options / Configure TeXstudio / Shortcuts} menüben.

\subsection{Kódrészletek}
Érdekesebb és bonyolultabb programozási megoldásainkat bátran illusztrálhatjuk kódrészletek beszúrásával. Fontos, hogy a beillesztett kódrészlet mérete álljon arányban annak fontosságával -- ritkán érdemes egy bekezdésnyi kódnál többet beszúrni egyszerre. Amennyiben egy bonyolultabb, akár több oldalas algoritmust szeretnénk bemutatni, annak kódját érdemesebb függelékbe rakni.

\subsubsection{Formázás}
A kódrészletek beillesztésére a \texttt{listings} csomagot használjuk (\texttt{lstlisting} környezet). A sablonban a csomag alapértelmezetten kulcsszókiemelést és monospace betűtípust alkalmaz. Kódrészletek beillesztésére semmiképpen ne képeket használjunk.

A kódrészletek formázásánál kerüljük a helypazarlást, illetve próbáljunk megelőzni az olvashatóságot rontó sortördelést, akár a forráskód módosításának árán is. Másolás előtt érdemes a behúzások mértékét 4-ről 2 karakterre csökkenteni.

Az alábbi C\# kódrészlet szemlélteti a \texttt{lstlisting} környezet használatát:

\begin{lstlisting}[language={[Sharp]C}]
static void Main(string[] args)
{
  var ci = new CultureInfo("en-us");
  ci.NumberFormat.CurrencySymbol = "";
  Thread.CurrentThread.CurrentCulture = ci;
  Console.WriteLine(ci);
  Console.WriteLine("{0:c}", 5.66);
}
\end{lstlisting}

A programozási elemek (osztálynevek, függvények) szövegközi kiemelésére használd a \verb|\texttt{NévMinták}| parancsot; ez egységes, monospace megjelenést biztosít.

\subsubsection{Irodalomjegyzék}
Az irodalomjegyzékben szereplő hivatkozásokat a \texttt{biblatex} csomag kezeli. A \verb|\printbibliography| parancs automatikusan generálja a listát a \texttt{references.bib} fájl tartalma alapján. A fordítási sorozatban a \texttt{biber} futtatása frissíti az irodalomjegyzéket. Csak azok a tételek jelennek meg a listában, amelyekre a szövegben \verb|\cite|-tal hivatkoztál -- a \verb|\nocite{*}| paranccsal az összes tétel megjeleníthető.

\subsection{Utolsó simítások}
Miután elkészültünk a dokumentációval, ne felejtsük el a következő lépéseket:
\begin{itemize}
  \item \textbf{Teljes fordítási sorozat:} \texttt{pdflatex} $\to$ \texttt{biber} $\to$ \texttt{pdflatex} $\to$ \texttt{pdflatex}. Ellenőrizd, hogy nem jelent-e meg valahol az „undefined reference" vagy „Citation ... undefined" figyelmeztetés a logban.
  \item \textbf{PDF metaadatok:} A \verb|\hypersetup| paranccsal a \texttt{pdfauthor}, \texttt{pdftitle} és \texttt{pdfkeywords} mezők kitölthetők -- ezek a PDF-fájl tulajdonságaiban jelennek meg.
  \item \textbf{Kinézet ellenőrzése PDF-ben:} A legjobb teszt a végén, ha a kész PDF-et átolvasod, különös tekintettel az oldalelrendezésre, az ábrák pozíciójára, a tartalomjegyzék oldalszámaira és az irodalomjegyzékre.
\end{itemize}

\end{document}
